\begin{table}[htbp]
\centering
\caption{Crop Substitution Event Study, Oregon Only (Illustrative)}
\label{tab:cropsub_eventstudy}
\begin{tabular}{lcccc}
\toprule
Event year (calendar year) & Border $\times$ Year coef. & SE (county-cl.) & $p$-value \\
\midrule
$-3$ (2012) & $+0.776$ & $0.425$ & $0.068$ \\
$+2$ (2017) & $+0.812$ & $0.430$ & $0.059$ \\
$+7$ (2022) & $+0.731$ & $0.393$ & $0.063$ \\
\bottomrule
\end{tabular}
\begin{minipage}{0.95\textwidth}
\vspace{0.3em}
\footnotesize
\textit{Notes:} Outcome is log(harvested cropland acres), \emph{levels}, not the change used in
Table~\ref{tab:cropsub_border70}. $n=108$ (36 Oregon counties $\times$ 3 NASS years), only 3 border
counties (Baker, Malheur, Wallowa) at every year, so treat as illustrative rather than well-powered.
Unlike the other three outcomes in this paper, Oregon is the only state with usable crop-substitution
data, so there is no cross-state market-age clock to stack on -- this event study uses Oregon's own 3
NASS waves (2012, 2017, 2022) relative to its 2015 retail launch. The coefficient is nearly identical
across all three years, \emph{including 2012, three years before Oregon legalized}: border counties simply
have more harvested cropland than interior counties in level terms throughout the entire period (Baker,
Malheur, and Wallowa are large wheat-growing counties in eastern Oregon), a fixed characteristic unrelated
to legalization timing. This is precisely why the main cross-sectional analysis
(Table~\ref{tab:cropsub_border70}) uses the \emph{change} in cropland rather than its level: differencing
removes this stable level gap, which is exactly what makes it not the right measure for a treatment effect.
\end{minipage}
\end{table}
